\subsection{RK4 implementado en CUDA}

Para establecer el punto de referencia, partimos de una implementación de RK4, que fue provista por el director del proyecto. Implementada usando CUDA, y librerías auxiliares como thrust, para uso de memoria entre device y host; nrrd, para la lectura del archivo de entrada.

El código realiza 3 tareas fundamentales:

\begin{itemize}
	\item Procesamiento del archivo de entrada, donde usa HDF5 para tomar los datos del campo y almacenarlos en un arreglo bidimensional, para optimizar la lectura en memoria.

	\item Kernel RK4, que consta de las operaciones relacionadas con el algoritmo RK4, que son realizadas directamente en el device, y sus resultados, almacenados en memoria del host

	\item Postprocesamiento para visualización, usa los datos de salida del kernel para producir un modelo 3D en formato VTK, de forma opcional, para que pueda ser visualizado.
\end{itemize}

Asimismo, el código tiene 3 parámetros de entrada, excluyendo el archivo H5, que serán utilizados posteriormente para la realización de pruebas:



\begin{table}[H]
	% \small
	\centering
	\vspace{-4mm}
	\caption{\\Parámetros de entrada del ejecutable} \label{tab:criterios}
	% \vspace{4mm}
	\begin{tabular}{lm{4.7cm}m{0.60\linewidth}c}
		\hline
		\multicolumn{1}{l}{}                   &
		\multicolumn{1}{c}{\textbf{Parámetro}} &
		\multicolumn{1}{c}{\textbf{Explicación}}                                                                                                                                      \\ \hline
		                                       & \centering num-seeds  &
		Corresponde al número de semillas generadas durante el paso inicial previo al kernel RK4                                                                                      \\
		                                       & \centering  num-steps &
		El número de veces o iteraciones de cálculos de pasos de RK4 por los que pasará cada semilla. Haciendo que el número total que se calcula RK4 es \textbf{num-seeds*num-steps} \\
		                                       & \centering vtk        &
		un valor numérico para que se guarde el archivo de visualización                                                                                                              \\
		\\ \hline
	\end{tabular}
	\vspace{-4mm}
\end{table}

No obstante, tras un análisis inicial del código, se optó por intercambiar la libería NRRD por HDF5, como sugerencia del director y autor del código, puesto que HDF5 posee características que se adecúan más a su uso en un contexto de computación heterogénea, como lectura/escritura paralela basada en MPI.

En adición, para la toma de métricas de tiempo de ejecución, se tomó el tiempo que le tomó ejecutar las dos primeras tareas del código, para medir aspectos como la lectura de archivos, escritura a memoria, tiempo de ejecución dentro del kernel y comunicación entre host y device. Tras este cambio, asumimos esta implementación como un punto de referencia, para derivar las implementaciones en los demás estándares a analizar.
